\chapter{表格排版}\label{ch:tables}
\section{默认设置与个性化设置}
\entry{设置 → 元素 → 表格设置… → 应用范围}
应用范围可选择“文档默认”或正文中的任一表格，无需先将光标放在表格内。文档默认影响沿用默认的表格及后续新增表格；个性化设置只作用于所选表格。确定后应用于所选范围，切换范围会重新载入该范围已保存的设置。

当前表格可以恢复沿用默认，取消独立设置。已有项目保存自己的尺寸，打开旧项目不会强制套用新默认值。

\section{两种表格源码}
简单表格用管道语法。第二行分隔表头和数据；冒号表示对齐。
\par\noindent\begin{minipage}{\linewidth}
\begin{Source}
| 1 | 2 | 3 |
| :--- | :---: | ---: |
| 12.50 | 13.25 | 14.00 |
| 15.00 | 16.75 | 17.50 |
\end{Source}
\end{minipage}\par\medskip
需要合并单元格时，使用显式闭合标签的 HTML 表格：
\par\noindent\begin{minipage}{\linewidth}
\begin{Source}
<table>
  <tr><th colspan="2">1</th><th>2</th></tr>
  <tr><td rowspan="2">3</td><td>4.5</td><td>6.7</td></tr>
  <tr><td>8.9</td><td>10.1</td></tr>
</table>
\end{Source}
\end{minipage}\par\medskip
\code{colspan} 表示跨列，\code{rowspan} 表示跨行。可使用受支持的单元格水平/垂直对齐和行内公式，不支持嵌套表格或任意 CSS。框线由排版器绘制，相邻公共边不会重复描线。

\section{默认值与列宽选择}
\par\noindent\begin{tabularx}{\linewidth}{@{}P{37mm}P{25mm}Y@{}}
\toprule
参数 & 默认值 & 使用建议\\\midrule
表格总宽 & 160 mm & 先核对可用正文宽度；窄模板需手动调小。\\
最小行高 & 8 mm & 内容可以撑高；启用固定行高后放不下就报错。\\
数字高度 & 3.8 mm & 独立于正文字高和公式字号。\\
四边留白 & 各 0.3 mm & 控制笔迹与框线净空，不建议为压缩表格降到重叠。\\
\bottomrule
\end{tabularx}

列宽可按内容分配、等宽、比例或固定毫米值。固定列宽的合计必须等于总宽。逐列对齐可填“沿用、左、中、右、小数点”，使用逗号分隔。小数点对齐适合小数位数不同的读数；它不会把 \code{001.20} 改写为 \code{1.2}。

当表格放不下时，先查是哪一列内容太长，再调整列宽、数字高度、公式字号或内容组织。程序不会静默压缩字形来填满单元格，也不会自动缩小整表。

\section{跨页与合并单元格}
普通行按可用页高分页，表头可重复。跨行合并连接的若干行必须作为一个完整行组同页排版，不会拆开单元格内部内容。如果整个行组高于完整正文区，继续增加页数也无济于事；需要拆分合并、拆表或调整正文区与字号。

重复表头属于原内容的重复展示，不是新输入的数据。检查分页时重点观察最后一行、续页首行以及长公式所在行，避免仅看第一页就导出。

